\documentclass{tufte-handout}

%\geometry{showframe}% for debugging purposes -- displays the margins

\usepackage{amsmath}

% Set up the images/graphics package
\usepackage{graphicx}
\setkeys{Gin}{width=\linewidth,totalheight=\textheight,keepaspectratio}
\graphicspath{{figures/}}

\title{Nitrogen Containing Compounds}
\author{Olivia Anderson and Dave Bridges}
%\date{}  % if the \date{} command is left out, the current date will be used

% The following package makes prettier tables.  We're all about the bling!
\usepackage{booktabs}

% The units package provides nice, non-stacked fractions and better spacing
% for units.
\usepackage{units}

% The fancyvrb package lets us customize the formatting of verbatim
% environments.  We use a slightly smaller font.
\usepackage{fancyvrb}
\fvset{fontsize=\normalsize}

% Small sections of multiple columns
\usepackage{multicol}

% Provides paragraphs of dummy text
\usepackage{lipsum}

% Nomenclature package
\usepackage[intoc]{nomencl}
\makenomenclature

% These commands are used to pretty-print LaTeX commands
\newcommand{\doccmd}[1]{\texttt{\textbackslash#1}}% command name -- adds backslash automatically
\newcommand{\docopt}[1]{\ensuremath{\langle}\textrm{\textit{#1}}\ensuremath{\rangle}}% optional command argument
\newcommand{\docarg}[1]{\textrm{\textit{#1}}}% (required) command argument
\newenvironment{docspec}{\begin{quote}\noindent}{\end{quote}}% command specification environment
\newcommand{\docenv}[1]{\textsf{#1}}% environment name
\newcommand{\docpkg}[1]{\texttt{#1}}% package name
\newcommand{\doccls}[1]{\texttt{#1}}% document class name
\newcommand{\docclsopt}[1]{\texttt{#1}}% document class option name

% Conditionally redefine Tufte-style macros for HTML output
\ifdefined\htmlversion
  \def\newthought#1{@@newthought:#1@@}
  \def\marginnote#1{@@marginnote:#1@@}
  \def\sidenote#1{@@sidenote:#1@@}
  \def\alttext#1{@@alttext:#1@@}
\fi
\providecommand{\alttext}[1]{}

\begin{document}

\maketitle% this prints the handout title, author, and date

\begin{abstract}
\noindent
Amino acids serve as scaffolds for a large collection of nitrogen-containing non-protein (NPN)\nomenclature{NPN}{Non-Protein Nitrogen} compounds; antioxidants, methyl donors, energy buffers, neurotransmitters, and other signaling molecules. This unit walks through the major NPN classes and, for each, asks a common set of questions: what is the pathway, what enzyme and cofactor is rate-limiting, what is the function, and how does it fail in practice? A frequent misconception is that failure is amino-acid-limited. In a modern diet it almost never is.
\end{abstract}

\tableofcontents
\pagebreak\index{Nitrogen-Containing Compounds|(}

\section{Learning Objectives}
\begin{itemize}
\item Understand the biosynthesis of nitrogen-containing non-protein compounds from amino acids
\item Explain the roles of glutathione, carnitine, choline, and creatine in the body
\item Identify the shared cofactor economies (SAM/one-carbon, BH4, NADPH, iron) that gate NPN biosynthesis
\item Apply a common four-mode framework (\textit {i.e.} cofactor limitation, genetic biosynthetic defect, demand state, drug/inhibitor interaction) to any NPN compound
\item Explain why precursor amino acid availability rarely rate-limits neurotransmitter synthesis, and describe the roles of BH4, MAO, COMT, MAOIs and SSRIs
\item Describe how heme biosynthesis and catabolism connect the protein and lipid units, and how the tryptophan-kynurenine route feeds NAD\textsuperscript{+} de novo synthesis
\end{itemize}

\section{A Framework for Understanding NPN Dysfunction}\label{sec:npn-framework}

\newthought{These compounds are rarely limited by the amino acid pool\sidenote{This is often surprising as supplements often conflate that precursors are necessary. For example while tyrosine is a precursor for dopamine, there is no evidence that dietary supplementation of tyrosine increases dopamine synthesis in the brain.  If you see any of these technically true but misleading \emph{structure--function} claims post them here}.} In a modern diet, the precursor amino acids for GSH, carnitine, choline, creatine, and neurotransmitters are essentially always available.\sidenote{The one important exception is \textbf{choline}, where population-level intake is genuinely below the Adequate Intake, especially in pregnancy. We will foreground this exception in its section.} When NPN pathways fail clinically, look for one of four causes:

\begin{enumerate}
\item \textbf{Cofactor limitation} --- especially SAM and the one-carbon economy (\textit{i.e.} folate, B12, methionine), BH4 (tetrahydrobiopterin) for the aromatic amino acid hydroxylases and NOS, NADPH for redox recycling, and iron/heme for oxygen-dependent enzymes.  The lack of cofactors can reduce the production of these compounds even when the amino acid precursors are abundant.
\item \textbf{Genetic biosynthetic defects} --- often the most common cause of a real clinical picture. Primary carnitine deficiency, GAMT/AGAT/CRTR defects in creatine metabolism, and G6PD deficiency limiting GSH regeneration are all more clinically relevant than dietary insufficiency of the corresponding amino acids.  Generally these can not be alleviated by precursor supplementation.
\item \textbf{Increased demand states} --- pregnancy, growth, oxidative stress, high-intensity exercise, sepsis, and aging can all outstrip otherwise adequate biosynthesis.
\item \textbf{Drug and inhibitor interactions} --- MAOIs and SSRIs on serotonin turnover, valproic acid on carnitine, statins on CoQ\textsubscript{10}, and acetaminophen on GSH are all clinically potent perturbations that are more clinically relevant than typical dietary effects.
\end{enumerate}

\newthought{Every section that follows uses this frame.} As you read about each compound, ask yourself which of the four modes actually explains its most common clinical failure.

\section{Review of Amino Acid Structure and Role in the Body}
An amino acid consists of four functional groups: a carboxyl, a hydrogen, an amino, and a side (R) group. The amino group makes the amino acid a nitrogen-containing structure. Amino acids are the building blocks of proteins, but they can also be used for the synthesis of NPN compounds either endogenously biosynthesized or, for some compounds, obtained from dietary sources. Once made, these compounds serve as antioxidants, methyl donors, transporters, neurotransmitters, or other signaling molecules. In the sections below we discuss glutathione, carnitine, choline, creatine, the major neurotransmitters, heme and polyamines, and the kynurenine route to NAD\textsuperscript{+}.

\section{Antioxidant Defense: Glutathione}
\newthought{Glutathione is synthesized in the liver} from glutamate\index{Amino Acids!Glutamate}, cysteine\index{Amino Acids!Cysteine}, and glycine\index{Amino Acids!Glycine}. It plays a central role as an intracellular antioxidant and is found in the majority of cells. The cysteine of glutathione\index{Metabolites!Glutathione} provides the structure with a thiol (-SH) group. The majority of cellular glutathione is in the reduced form, but it can be oxidized in the presence of free radicals and lipid peroxides. Glutathione donates its hydrogen to stabilize these species; following oxidation, NADPH\index{Cofactors!NADPH} from the pentose phosphate pathway\index{Pentose Phosphate Pathway}\sidenote{See the pentose phosphate pathway chapter for the detailed treatment of the GSH/GSSG cycle and its dependence on NADPH.} rapidly reduces glutathione back. The ratio of reduced glutathione (GSH)\nomenclature{GSH}{Glutathione (Reduced)} to oxidized glutathione (GSSG)\nomenclature{GSSG}{Glutathione (Oxidized)} is a widely used biomarker of oxidative stress and inflammation\index{Inflammation}. In hepatocytes under stress or damage, tissue GSH:GSSG falls \citep{zitkaRedoxStatusExpressed2012,sentellasGSSGGSHRatios2014}.

\newthought{Where does GSH biology break?} Applying the four-mode frame:

\begin{itemize}
\item \emph{Cofactor limitation.} Cysteine availability, not glutamate or glycine, sets the ceiling on synthesis. Cysteine is conditionally essential and depends on methionine\index{Amino Acids!Methionine} flux through the transsulfuration pathway.\sidenote{This is why N-acetylcysteine (NAC)\nomenclature{NAC}{N-Acetylcysteine} is used therapeutically: it delivers cysteine directly. In acute acetaminophen\index{Drugs!Acetaminophen} overdose, NAC is the standard antidote precisely because it restores GSH synthesis capacity.} NADPH availability sets the ceiling on GSH \emph{regeneration}.
\item \emph{Genetic defects.} G6PD deficiency\index{Diseases!G6PD Deficiency}\nomenclature{G6PD}{Glucose-6-Phosphate Dehydrogenase} limits NADPH generation and therefore GSH regeneration; glutathione synthetase (GSS)\nomenclature{GSS}{Glutathione Synthetase}\index{Enzymes!Glutathione Synthetase} mutations are rarer but produce severe presentations.
\item \emph{Demand.} Chronic oxidative stress\index{Oxidative Stress}, inflammation, and rapid proliferation deplete the GSH pool.
\item \emph{Drug/inhibitor.} Acetaminophen at toxic doses is metabolized to NAPQI\nomenclature{NAPQI}{N-Acetyl-p-Benzoquinone Imine}, which is conjugated to GSH; overdose depletes hepatic GSH and produces the classic centrilobular necrosis. NAC rescue restores cysteine supply and re-enables GSH synthesis.
\end{itemize}

\section{The Methylation Economy}\label{sec:methylation-economy}
\newthought{Carnitine, choline, and creatine share a hidden dependency:} their biosynthesis all consumes methyl groups from S-adenosyl methionine (SAM)\nomenclature{SAM}{S-Adenosyl Methionine}\index{Cofactors!SAM}\index{S-Adenosyl Methionine} via one-carbon metabolism\index{One Carbon Metabolism}. Failure of the methylation economy, through choline, folate or B12 insufficiency, methionine restriction, or common polymorphisms in one-carbon enzymes, can appear as coordinated deficits in all three, and it can also elevate homocysteine\index{Metabolites!Homocysteine}.

\subsection{Carnitine}
Carnitine is synthesized in the liver and, to a lesser extent, the kidneys. Three methyl groups are added to lysine\index{Amino Acids!Lysine} by SAM via one-carbon metabolism, and the methylated lysine is then hydroxylated to form carnitine. Once synthesized, carnitine is stored predominantly in muscle, which is why we can also obtain it from dietary animal sources (meat, fish). Approximately 60--80\% of ingested carnitine is absorbed at typical intake levels of 0.5--0.6~g \citep{reboucheKineticsPharmacokineticsRegulation2004}. Carnitine\index{Carnitine} is available in supplemental form.

\newthought{Carnitine's main role} is to transport long-chain fatty acids across the inner mitochondrial membrane. The fatty acid is covalently joined to carnitine via carnitine palmitoyltransferase I (CPT-I)\sidenote{Sometimes referred to as CAT (carnitine acyltransferase).}\nomenclature{CPT-I}{Carnitine Palmitoyltransferase I}; once inside, CPT-II\nomenclature{CPT-II}{Carnitine Palmitoyltransferase II} releases the fatty acid for $\beta$-oxidation. This process is discussed in detail in the lipid oxidation chapter.

\newthought{Failure modes for carnitine:}
\begin{itemize}
\item \emph{Genetic.} Primary carnitine deficiency (SLC22A5/OCTN2)\index{Diseases!Primary Carnitine Deficiency}\nomenclature{OCTN2}{Organic Cation/Carnitine Transporter Novel 2 (SLC22A5)} impairs cellular uptake and causes cardiomyopathy, hypoketotic hypoglycemia, and muscle weakness. Long-chain fatty acid oxidation defects (\textit{e.g.}, VLCAD\nomenclature{VLCAD}{Very Long-Chain Acyl-CoA Dehydrogenase}, LCHAD\nomenclature{LCHAD}{Long-Chain 3-Hydroxyacyl-CoA Dehydrogenase}) produce a secondary picture \citep{longoCarnitineTransportFatty2016}.
\item \emph{Drug.} Valproic acid\index{Drugs!Valproic Acid} depletes carnitine through conjugation and is a well-recognized cause of secondary deficiency.
\item \emph{Demand.} Neonates on prolonged carnitine-free parenteral nutrition are the only common setting where dietary intake truly limits carnitine.
\item \emph{Cofactor.} Because carnitine biosynthesis depends on SAM, folate or B12 insufficiency reduces synthesis capacity, though this is rarely the presenting problem.
\end{itemize}
Dietary insufficiency in healthy adults (including people eating vegetarian and vegan diets) reduces plasma carnitine modestly but is not clinically meaningful.\sidenote{Muscle stores respond to supplementation but performance benefits in trained individuals are inconsistent.}

\newthought{From here on, apply the framework yourself.}\sidenote{As you read the remainder of this chapter, use the four-mode framework of Section~\ref{sec:npn-framework} as a self-study exercise. For each of the remaining NPN compounds (choline, creatine, catecholamines, serotonin, histamine, GABA, acetylcholine, nitric oxide, heme, polyamines, and the kynurenine route to NAD\textsuperscript{+}), identify which failure mode is most clinically relevant, and what specific gene, drug, or demand state anchors it. The prose that follows describes the biology; the four-mode diagnosis is left to you.} The two worked examples above (glutathione and carnitine) are the pattern. In every remaining section below the biology is described, but the failure-mode analysis is left as an exercise.

\subsection{Choline: the Notable Exception}
Choline is synthesized in the liver and, to a lesser extent, the kidneys. The endogenous route begins with phosphatidylserine\index{Phospholipids!Phosphatidylserine}, which is decarboxylated to phosphatidylethanolamine (PE)\nomenclature{PE}{Phosphatidylethanolamine}\index{Phospholipids!Phosphatidylethanolamine}. Three methyl groups from SAM are then transferred to PE by phosphatidylethanolamine N-methyltransferase (PEMT)\nomenclature{PEMT}{Phosphatidylethanolamine N-Methyltransferase}\index{Enzymes!PEMT}, yielding phosphatidylcholine\index{Choline}\index{Phospholipids!Phosphatidylcholine}; the principal internal source of choline. Choline is also abundant in dietary sources such as eggs, wheat germ, legumes, salmon, and liver, either as free choline or bound in lecithin.

\newthought{Choline is the exception to the ``diet doesn't matter'' rule.} Most Americans consume less than the Adequate Intake \citep{wallaceUsualCholineIntakes2017}, and choline demand rises further in pregnancy and lactation. Common PEMT polymorphisms shift individual dependence on dietary intake. Because phosphatidylcholine is required for VLDL assembly, choline insufficiency directly impairs hepatic export of triglyceride, producing hepatic steatosis a cause of MASLD\nomenclature{MASLD}{Metabolic Dysfunction-Associated Steatotic Liver Disease}\index{Diseases!MASLD}.\sidenote{MASLD was previously known as NAFLD. Choline is one of the few dietary factors with a mechanistically clean and reproducible link to hepatic steatosis in humans.}

\paragraph{Choline in One-Carbon Metabolism}
Choline is not only a product of SAM-dependent methylation, it is also, via betaine\index{Metabolites!Betaine}, a methyl \emph{donor} that feeds one-carbon metabolism. Adequate choline is therefore required for one-carbon metabolism to run efficiently. Poor choline status contributes to elevated blood homocysteine\sidenote{Distinguish \emph{hyperhomocysteinemia} (a modest elevation associated with B-vitamin insufficiency and increased atherosclerotic risk) from \emph{homocystinuria}\index{Diseases!Homocystinuria}, the specific inborn error of cystathionine $\beta$-synthase (CBS)\nomenclature{CBS}{Cystathionine $\beta$-Synthase}\index{Enzymes!CBS}. They are different clinical entities driven by different mechanisms.}, which is associated with vascular wall damage, atherosclerosis, stroke, and heart disease \citep{millardDietaryCholineBetaine2018}. Reduced flux through one-carbon metabolism can also alter nucleic acid synthesis and cell division, contributing to increased risk of neural tube defects\index{Diseases!Neural Tube Defects} such as spina bifida \citep{zeiselCholineCriticalRole2006}.

\paragraph{Choline as Structural Lipid: MASLD}
Phosphatidylcholine (PC) is the dominant phospholipid on the surface of very low-density lipoprotein (VLDL)\nomenclature{VLDL}{Very Low-Density Lipoprotein}\index{Lipoproteins!VLDL} particles. When choline is insufficient, hepatic PC synthesis falls, VLDL assembly is impaired, and triglyceride accumulates in the liver. This is the mechanistic basis for choline-deficient models of hepatic steatosis and remains one of the clearest dietary contributors to MASLD in humans.

\newthought{Individual choline requirements vary widely.}  Common variants in \emph{PEMT}, \emph{MTHFD1}\nomenclature{MTHFD1}{Methylenetetrahydrofolate Dehydrogenase 1}\index{Enzymes!MTHFD1}, and \emph{CHDH}\nomenclature{CHDH}{Choline Dehydrogenase}\index{Enzymes!CHDH} raise the individual dietary requirement, particularly for women, and pregnancy and lactation raise requirements substantially maternal deficiency has offspring consequences \citep{craciunescuDietaryCholineReverses2010}. Chronic parenteral nutrition without choline was historically a major cause of iatrogenic deficiency.

\subsection{Creatine}
In the kidney, glycine and arginine\index{Amino Acids!Arginine} condense to form guanidinoacetate\index{Metabolites!Guanidinoacetate}. Guanidinoacetate is then methylated by SAM in the liver via guanidinoacetate N-methyltransferase (GAMT)\nomenclature{GAMT}{Guanidinoacetate N-Methyltransferase}\index{Enzymes!GAMT} to yield creatine. Creatine is then transported to muscle, where about 95\% of the body pool resides. Dietary sources are almost exclusively meat and fish, and creatine is widely used as a performance-enhancement supplement.

\newthought{Creatine's role} is to act as a high-energy phosphate buffer via phosphocreatine\index{Metabolites!Phosphocreatine}. More than half of the creatine in resting muscle is phosphorylated. During contraction, phosphocreatine transfers its phosphate to ADP, generating ATP quickly and delaying the onset of glycogenolysis. Between contractions, creatine kinase\index{Enzymes!Creatine Kinase} runs the reverse direction to reload phosphocreatine, so the reaction is bidirectional depending on the ATP/ADP ratio.\sidenote{This phosphocreatine buffer supports maximal effort for roughly 8--30 seconds, which is why creatine supplementation reliably improves short-duration, high-intensity performance (sprinting, weight lifting; but not necessarily endurance.} Creatine and phosphocreatine are slowly and non-enzymatically converted to creatinine\index{Metabolites!Creatinine}, which is transferred from muscle to the kidney and excreted in urine. Urinary creatinine is used as a biomarker of muscle mass and kidney function \citep{kreiderInternationalSocietySports2017}.

\newthought{Three creatine deficiency syndromes} illustrate the genetic axis: AGAT, GAMT, and creatine transporter (CRTR/SLC6A8) deficiency\index{Diseases!Creatine Deficiency Syndromes}\nomenclature{AGAT}{Arginine:Glycine Amidinotransferase}\nomenclature{CRTR}{Creatine Transporter (SLC6A8)} all present with intellectual disability, and AGAT/GAMT respond to oral creatine.\sidenote{CRTR deficiency does not respond to supplementation because the transporter defect blocks brain uptake. These are the modal clinically meaningful ``creatine deficiencies.''} Vegans and vegetarians show lower plasma and muscle creatine but functional consequences are limited and respond briskly to supplementation.

\subsection{When the Methylation Supply System Breaks Down}
\newthought{A shared failure mode across carnitine, choline, and creatine} is inadequate SAM/one-carbon flux. Vitamin B12 deficiency\index{Vitamins!Vitamin B12 (Cobalamin)} (from pernicious anemia, chronic PPI\nomenclature{PPI}{Proton Pump Inhibitor}\index{Drugs!PPI} use, or dietary avoidance) traps folate as methyl-tetrahydrofolate, reducing methionine regeneration and therefore SAM. Folate\index{Vitamins!Vitamin B9 (Folate)} deficiency has the same downstream effect through a different route. Methionine-restricted diets amplify the problem. Common polymorphisms in \emph{MTHFR}\nomenclature{MTHFR}{Methylenetetrahydrofolate Reductase}\index{Enzymes!MTHFR}, \emph{PEMT}, and related genes shift individual sensitivity. In practice, biomarker patterns can look coordinated: elevated homocysteine, low plasma choline, low carnitine, and lower creatine synthesis all co-occur. This shared cofactor economy is more clinically useful to understand than any of these compounds in isolation.

\section{Signaling Molecules}\label{sec:signaling}

\subsection{Tetrahydrobiopterin (BH4): The Underappreciated Cofactor}
\newthought{Before discussing individual neurotransmitters, a note on BH4.} Tetrahydrobiopterin\index{Cofactors!Tetrahydrobiopterin (BH4)}\nomenclature{BH4}{Tetrahydrobiopterin} is the obligate cofactor for the aromatic amino acid hydroxylases (phenylalanine hydroxylase (PAH)\nomenclature{PAH}{Phenylalanine Hydroxylase}\index{Enzymes!PAH}, tyrosine hydroxylase (TH), and tryptophan hydroxylase (TPH)) and for all three nitric oxide synthase isoforms (nNOS, eNOS, iNOS). BH4 availability, not precursor amino acid availability, is often the true rate-limiter for catecholamine and serotonin synthesis, and for endothelial NO production. BH4 deficiency\index{Diseases!BH4 Deficiency} can cause a variant of hyperphenylalaninemia, and BH4 is now a therapeutic (sapropterin) for a subset of PKU\nomenclature{PKU}{Phenylketonuria}\index{Diseases!PKU} patients. BH4 oxidation under oxidative stress uncouples eNOS and is thought to contribute to endothelial dysfunction\index{Diseases!Endothelial Dysfunction} making a bridge between redox biology (the GSH story above) and vascular disease.

\subsection{Catecholamines: Tyrosine \texorpdfstring{$\rightarrow$}{->} Dopamine \texorpdfstring{$\rightarrow$}{->} Norepinephrine \texorpdfstring{$\rightarrow$}{->} Epinephrine}
\newthought{Catecholamine biosynthesis} begins with tyrosine\index{Amino Acids!Tyrosine}.\sidenote{Tyrosine itself is derived from phenylalanine\index{Amino Acids!Phenylalanine} by PAH, another BH4-dependent step. This is why classical PKU produces both hyperphenylalaninemia and secondary catecholamine and pigmentation defects.} Tyrosine hydroxylase (TH)\index{Enzymes!Tyrosine Hydroxylase}\nomenclature{TH}{Tyrosine Hydroxylase}, requiring BH4 and iron\index{Minerals!Iron}, converts tyrosine to L-DOPA\index{Metabolites!L-DOPA}. Aromatic L-amino acid decarboxylase (AADC)\nomenclature{AADC}{Aromatic L-Amino Acid Decarboxylase}\index{Enzymes!AADC}, requiring pyridoxal 5'-phosphate (PLP)\nomenclature{PLP}{Pyridoxal 5'-Phosphate}\index{Cofactors!Pyridoxal Phosphate (PLP)}\index{Vitamins!Vitamin B6 (Pyridoxine)}, then produces dopamine\index{Neurotransmitters!Dopamine}. Dopamine $\beta$-hydroxylase (DBH)\nomenclature{DBH}{Dopamine $\beta$-Hydroxylase}\index{Enzymes!DBH} converts dopamine to norepinephrine\index{Neurotransmitters!Norepinephrine}, and phenylethanolamine N-methyltransferase (PNMT)\nomenclature{PNMT}{Phenylethanolamine N-Methyltransferase}\index{Enzymes!PNMT} methylates norepinephrine to epinephrine\index{Neurotransmitters!Epinephrine} using SAM. TH is the rate-limiting enzyme and is under end-product feedback inhibition. Precursor tyrosine availability is rarely limiting in a normal diet.

\subsection{Serotonin and Melatonin: Tryptophan \texorpdfstring{$\rightarrow$}{->} 5-HTP \texorpdfstring{$\rightarrow$}{->} Serotonin \texorpdfstring{$\rightarrow$}{->} Melatonin}
\newthought{Serotonin biosynthesis} parallels the catecholamine route. Tryptophan hydroxylase (TPH)\index{Enzymes!Tryptophan Hydroxylase}\nomenclature{TPH}{Tryptophan Hydroxylase}, again BH4- and iron-dependent, produces 5-hydroxytryptophan (5-HTP)\nomenclature{5-HTP}{5-Hydroxytryptophan}\index{Metabolites!5-HTP} from tryptophan\index{Amino Acids!Tryptophan}, which is decarboxylated by AADC (PLP) to serotonin (5-HT)\nomenclature{5-HT}{5-Hydroxytryptamine (Serotonin)}\index{Neurotransmitters!Serotonin}. In the pineal gland, serotonin is subsequently N-acetylated and O-methylated to melatonin\index{Neurotransmitters!Melatonin}, providing the circadian signal.

\newthought{Serotonin turnover, MAO, and pharmacology.} Serotonin (and other monoamines) are oxidatively deaminated by monoamine oxidase (MAO)\index{Enzymes!MAO}\nomenclature{MAO}{Monoamine Oxidase} to their aldehyde intermediates, subsequently converted to 5-HIAA\nomenclature{5-HIAA}{5-Hydroxyindoleacetic Acid}\index{Metabolites!5-HIAA} and excreted. Catechol-O-methyltransferase (COMT)\index{Enzymes!COMT}\nomenclature{COMT}{Catechol-O-Methyltransferase} contributes to catecholamine catabolism using SAM as the methyl donor. MAO inhibitors (MAOIs)\nomenclature{MAOI}{Monoamine Oxidase Inhibitor}\index{Drugs!MAOI} are an older class of antidepressant that raises synaptic serotonin (and norepinephrine and dopamine) by blocking oxidative deamination. Patients taking MAOIs must eliminate tyramine\index{Metabolites!Tyramine}-rich foods\sidenote{Aged cheeses, cured meats, fermented soy, some red wines and draft beers.}, because dietary tyramine is normally destroyed by intestinal MAO. If systemic MAO is blocked, dietary tyramine reaches the circulation, displaces norepinephrine from sympathetic nerve terminals, and can precipitate a hypertensive crisis. Selective serotonin reuptake inhibitors (SSRIs)\nomenclature{SSRI}{Selective Serotonin Reuptake Inhibitor}\index{Drugs!SSRI}, a much more commonly used class today, act instead at the presynaptic serotonin transporter (SERT)\nomenclature{SERT}{Serotonin Transporter}\index{Enzymes!SERT}.

\subsection{Histamine, GABA, and Acetylcholine}
\newthought{Histamine\index{Neurotransmitters!Histamine}} is synthesized from histidine\index{Amino Acids!Histidine} by histidine decarboxylase (HDC)\nomenclature{HDC}{Histidine Decarboxylase}\index{Enzymes!HDC}, PLP-dependent, and mediates gastric acid secretion (via H2 receptors on parietal cells), inflammatory responses, and neurotransmission.

\newthought{GABA\index{Neurotransmitters!GABA}}\nomenclature{GABA}{$\gamma$-Aminobutyric Acid}, the principal inhibitory neurotransmitter in the CNS, is made from glutamate by glutamate decarboxylase (GAD)\nomenclature{GAD}{Glutamate Decarboxylase}\index{Enzymes!GAD}, also PLP-dependent. This is one reason why B6\index{Vitamins!Vitamin B6 (Pyridoxine)} deficiency produces seizures.

\newthought{Acetylcholine\index{Neurotransmitters!Acetylcholine}} synthesis takes free choline (see above) and transfers an acetyl group from acetyl-CoA using choline acetyltransferase (ChAT)\nomenclature{ChAT}{Choline Acetyltransferase}\index{Enzymes!Choline Acetyltransferase}. Acetylcholine is central to autonomic control, neuromuscular transmission, and cognitive functions including long-term memory. Choline deficiency during pregnancy has been associated in animal studies with structural brain abnormalities and impaired long-term memory in offspring \citep{craciunescuDietaryCholineReverses2010}.

\subsection{Nitric Oxide from Arginine}
Nitric oxide (NO)\index{Signaling Molecules!Nitric Oxide}\nomenclature{NO}{Nitric Oxide} is generated from arginine by nitric oxide synthases nNOS\nomenclature{nNOS}{Neuronal Nitric Oxide Synthase}, eNOS\nomenclature{eNOS}{Endothelial Nitric Oxide Synthase}, and iNOS\nomenclature{iNOS}{Inducible Nitric Oxide Synthase}\index{Enzymes!Nitric Oxide Synthase}\nomenclature{NOS}{Nitric Oxide Synthase (generic)} all BH4- and NADPH-dependent. Endothelial NO from eNOS is a key vasodilator and platelet inhibitor. Under oxidative stress, oxidized BH4 uncouples eNOS so that it produces superoxide instead of NO; linking the GSH/NADPH story from Chapter~\ref{sec:npn-framework} to endothelial dysfunction and atherosclerosis.

\section{Heme, Porphyrins, and Polyamines}\label{sec:heme-poly}

\subsection{Heme Biosynthesis}
\newthought{Heme biosynthesis} begins in the mitochondrion. Succinyl-CoA\index{Metabolites!Succinyl-CoA} (from the TCA cycle) and glycine condense to form $\delta$-aminolevulinate (ALA)\nomenclature{ALA}{$\delta$-Aminolevulinate}\index{Metabolites!ALA} via ALA synthase (ALAS)\nomenclature{ALAS}{$\delta$-Aminolevulinate Synthase}\index{Enzymes!ALA Synthase}, the rate-limiting enzyme. ALAS is repressed by heme in a classic feedback loop. Two ALA molecules combine to form porphobilinogen\index{Metabolites!Porphobilinogen}, and successive condensations yield uroporphyrinogen and coproporphyrinogen, culminating in protoporphyrin IX\index{Metabolites!Protoporphyrin IX}. In the last step, iron is inserted by ferrochelatase\index{Enzymes!Ferrochelatase} to yield heme\index{Metabolites!Heme}. About 85\% of heme synthesis occurs in erythroid precursors for hemoglobin; the remainder supports cytochromes, catalase, and other hemoproteins.

\subsection{Porphyrias (brief)}
\newthought{The porphyrias\index{Diseases!Porphyrias}} are a family of inherited (or, in one case, acquired) defects in the enzymes of heme biosynthesis. Different enzymatic blocks produce accumulation of different intermediates and correspondingly different clinical pictures --- acute neurovisceral attacks (acute intermittent porphyria) or photosensitivity (porphyria cutanea tarda). The unifying teaching point is that impaired heme synthesis releases ALAS from feedback inhibition, driving overproduction of upstream intermediates \citep{puyPorphyrias2010}.

\subsection{Heme Catabolism and Bilirubin}
Old red blood cells\index{Red Blood Cells} are cleared in the reticuloendothelial system. Heme is opened by heme oxygenase\index{Enzymes!Heme Oxygenase} to biliverdin\index{Metabolites!Biliverdin} (releasing iron for reuse), and biliverdin is reduced to unconjugated bilirubin\index{Metabolites!Bilirubin}. Bilirubin is transported bound to albumin to the liver, conjugated with glucuronic acid, and secreted in bile. This is where the protein and lipid units connect to gastrointestinal physiology.

\subsection{Polyamines}
\newthought{Polyamines\index{Metabolites!Polyamines}} --- putrescine\index{Metabolites!Putrescine}, spermidine\index{Metabolites!Spermidine}, and spermine\index{Metabolites!Spermine} are synthesized from ornithine\index{Amino Acids!Ornithine} (via arginase from arginine). Ornithine decarboxylase (ODC)\nomenclature{ODC}{Ornithine Decarboxylase}\index{Enzymes!ODC}, PLP-dependent, is the rate-limiting enzyme and has one of the shortest protein half-lives in mammalian cells. ODC is strongly induced during rapid proliferation (\textit{e.g.} growth, wound healing, cancer) and its inhibitor $\alpha$-difluoromethylornithine (DFMO)\nomenclature{DFMO}{$\alpha$-Difluoromethylornithine}\index{Drugs!DFMO} is used both as a therapeutic and as a chemoprevention agent under investigation. Polyamines stabilize nucleic acids and are required for translation, so their availability tracks growth demand rather than dietary supply.

\section{Tryptophan, Kynurenine, and de novo NAD\textsuperscript{+}}\label{sec:kynurenine-nad}
\newthought{Tryptophan is the exception} among the amino acid $\rightarrow$ NPN pathways because it yields NAD\textsuperscript{+}\index{Metabolites!NAD} via de novo synthesis. About 5\% of dietary tryptophan is diverted through the kynurenine\index{Metabolites!Kynurenine} pathway to nicotinic acid mononucleotide (NaMN)\nomenclature{NaMN}{Nicotinic Acid Mononucleotide}\index{Metabolites!NaMN}, which enters the NAD\textsuperscript{+} pool. The efficiency of this pathway is such that dietary niacin (vitamin B3)\index{Vitamins!Vitamin B3 (Niacin)} requirements can be partially met by adequate tryptophan intake\sidenote{This is why corn-based diets historically produced pellagra\index{Diseases!Pellagra} even without frank tryptophan deficiency (corn's tryptophan is bound and its niacin content is low).}. See the TCA cycle chapter for the modern discussion of NAD\textsuperscript{+} pool dynamics, sirtuins, and aging \citep{covarrubiasNADMetabolismIts2021}.

\section{Reflection Questions}
\begin{enumerate}
\item A patient presents to the emergency department with acute liver failure following an acetaminophen overdose. Using the four failure modes described in Section~\ref{sec:npn-framework}, explain the mechanism of liver damage in terms of the glutathione cycle, describe the biochemical rationale for N-acetylcysteine (NAC) rescue, and predict which co-existing condition (G6PD deficiency, alcoholism, chronic malnutrition) would be expected to worsen the prognosis and why.
\item A middle-aged woman on a low-egg, low-red-meat, high-fiber diet is diagnosed with MASLD despite a normal body weight and no diabetes. Using your knowledge of the choline biosynthesis pathway (PEMT), one-carbon metabolism, and VLDL assembly, explain the mechanistic connection between her dietary choices, her PEMT allele frequency, and the accumulation of hepatic triglyceride. Include a discussion of why supplementing folate or B12 might partially, but not fully, address the problem.
\item A vegan endurance athlete asks whether creatine supplementation would improve their performance in a marathon and in a subsequent 400-meter sprint. Distinguish the roles of the phosphocreatine energy buffer and aerobic ATP production, evaluate the likely magnitude of benefit for each event, and comment on whether their vegan diet would meaningfully limit muscle creatine stores at baseline.
\item A patient started on an MAOI for treatment-resistant depression develops a severe headache and blood pressure of 210/120 after a dinner of aged cheese and cured meats. Explain the mechanism using the biology of MAO in the intestinal wall, tyramine, and norepinephrine release from sympathetic nerve terminals. Contrast this with the mechanism of action of an SSRI and explain why the SSRI is not subject to the same food interaction.
\end{enumerate}

\section{Protein and Nitrogen Unit Integration Questions}\label{sec:protein-integration}
\begin{enumerate}
\item Trace a single leucine molecule from a bolus of whey protein through digestion (brush-border peptidases, transporter selection), absorption (BCAA-specific transporter kinetics), incorporation into muscle protein synthesis (mTORC1 activation), possible catabolism via BCKDH into acetyl-CoA and acetoacetate, and its potential contribution to ketogenesis. Compare this to the fate of a single tryptophan molecule from the same meal into serotonin, kynurenine/NAD\textsuperscript{+}, or full oxidation and use the two trajectories to illustrate why ``amino acids for protein synthesis'' is only one of several fates.
\item A patient with a urea cycle disorder (specifically, ornithine transcarbamylase deficiency) presents with hyperammonemia after a high-protein meal. Explain the block, predict the downstream effects on: (a) polyamine biosynthesis, (b) NO production, and (c) creatine biosynthesis. Then propose a dietary and pharmacological strategy that both reduces nitrogen load and preserves the NPN compounds these patients still need.
\item A vegan patient of Northern European ancestry presents with elevated homocysteine, low plasma choline, and modestly reduced muscle creatine. Using your knowledge of the shared methylation economy, propose a mechanistic explanation, identify the specific vitamin(s) and metabolite(s) most likely to be limiting, and design a rational dietary intervention. Address explicitly why supplementing methionine alone would not fully resolve the picture.
\item A sarcopenic 78-year-old on a proton pump inhibitor for reflux presents with progressive weakness, macrocytic anemia, and mild cognitive impairment. Connect the B12 malabsorption story (from the digestive chapter) to failure of the methylation economy (this chapter) and to the mitochondrial-biogenesis / NAD\textsuperscript{+} discussion (TCA cycle chapter). Propose which single intervention would have the highest probability of measurable clinical benefit and explain the mechanism.
\item An athlete with G6PD deficiency asks whether they should worry about intense exercise, high-dose vitamin C, and any supplement marketed as an ``antioxidant booster.'' Integrate the pentose phosphate pathway, GSH cycling, cysteine/methionine supply, and the endothelial NO/BH4 uncoupling story to advise them. Distinguish what would be truly risky from what is only theoretical concern.
\item Design an idealized nutrition protocol for the third trimester of pregnancy that specifically addresses the elevated demands for: (a) choline (fetal brain development, PEMT-limited PC supply), (b) folate/B12 (nucleotide synthesis, methylation economy), (c) iron (heme synthesis for both maternal and fetal erythropoiesis), and (d) BH4/PLP-dependent neurotransmitter synthesis. Explain the mechanistic basis for each recommendation, and identify which single nutrient's insufficiency would have the largest downstream cascade of consequences across the NPN compound families discussed in this chapter.
\end{enumerate}

\index{Nitrogen-Containing Compounds|)}

\bibliography{library}
\bibliographystyle{plainnat}

\end{document}
